\chapter*{Remerciements}
\markleft{Remerciements}
\addcontentsline{toc}{chapter}{Remerciements}

Christophe Rieunier a fait la route depuis Nantes pour nous parler d'analyse de
\emph{malware}. Il aurait pu dérouler son support et repartir. Au lieu de ça, il
est resté le soir pour nous montrer les conférences qu'il avait données au fil de
sa carrière, et pour raconter ce qu'il y avait vu. On est restés aussi. Pas parce
qu'on suivait tout, loin de là, mais parce qu'on n'avait pas envie que ça
s'arrête. Il y a des gens qui vous transmettent leur passion sans jamais chercher
à impressionner : il est de ceux-là. Si ce mémoire parle de \emph{malwares}, c'est
grâce à vous. Alors merci.

Merci à Reza El Galai, responsable pédagogique du Mastère spécialisé, d'avoir tenu
la promo à distance de Troyes sans jamais nous lâcher. Encadrer des gens qu'on ne
croise pas dans un couloir demande un effort qu'on ne remarque que le jour où il
n'est pas fait. Et c'est lui qui est allé chercher des intervenants comme
Christophe : un cours vaut d'abord par celui qui le donne, et ces choix-là ne se
font pas tout seuls.

Merci à Jean-Christophe Pillet de m'avoir pris comme alternant chez Login
Sécurité. Faire confiance à quelqu'un qui débute, ça coûte du temps, et il en a
donné sans compter, disponible pour mes questions, y compris celles dont je savais
en les posant qu'elles étaient bêtes. J'ai autant appris en travaillant dans son
équipe qu'en suivant des cours. C'est une chance que j'aimerais avoir la lucidité
de mesurer encore dans dix ans.

Merci enfin à ma famille, qui ne m'a pas regardé de travers quand j'ai annoncé que
je repartais en études. Et merci surtout à ma compagne, Mathilde, qui a tenu trois ans de
relation à distance et qui a relu ces pages. Trente pages sur la détection de
\emph{malware} quand on n'est pas du métier, c'est une preuve d'amour. Elle a tenu
bon, donc moi aussi.